% Generated cross-document labels; rebuild with build.py after edits.
\newlabel{S-sec:extended_selection}{{S2}{1}{Extended theory for certified finite-portfolio selection}{section.2}{PhysicaA_Supplement.pdf}}
\newlabel{S-sec:endpoint_selection}{{S2.1}{1}{Endpoint screening for a finite action portfolio}{subsection.2.1}{PhysicaA_Supplement.pdf}}
\newlabel{S-eq:portfolio_diameter}{{S1}{1}{Endpoint screening for a finite action portfolio}{equation.2.1}{PhysicaA_Supplement.pdf}}
\newlabel{S-eq:envelope_bet}{{S2}{1}{Endpoint screening for a finite action portfolio}{equation.2.2}{PhysicaA_Supplement.pdf}}
\newlabel{S-thm:endpoint_certificate}{{S1}{2}{Anytime endpoint certificate}{theorem.1}{PhysicaA_Supplement.pdf}}
\newlabel{S-eq:paired_bets}{{S3}{2}{A paired fallback with explicit unresolved outcomes}{equation.2.3}{PhysicaA_Supplement.pdf}}
\newlabel{S-thm:envelope_selection}{{S2}{2}{Envelope-first finite-portfolio selection}{theorem.2}{PhysicaA_Supplement.pdf}}
\newlabel{S-eq:portfolio_pac}{{S4}{2}{Envelope-first finite-portfolio selection}{equation.2.4}{PhysicaA_Supplement.pdf}}
\newlabel{S-alg:envelope_first}{{S1}{3}{A paired fallback with explicit unresolved outcomes}{algorithm.S1}{PhysicaA_Supplement.pdf}}
\newlabel{S-prop:zero_width_cost}{{S3}{3}{Zero-width endpoint cost}{theorem.3}{PhysicaA_Supplement.pdf}}
\newlabel{S-eq:zero_width_cost}{{S5}{3}{Zero-width endpoint cost}{equation.2.5}{PhysicaA_Supplement.pdf}}
\newlabel{S-sec:screen_cost}{{S2.3}{4}{When can endpoint screening repay its cost?}{subsection.2.3}{PhysicaA_Supplement.pdf}}
\newlabel{S-eq:chord_drift}{{S6}{4}{When can endpoint screening repay its cost?}{equation.2.6}{PhysicaA_Supplement.pdf}}
\newlabel{S-thm:screen_power}{{S5}{4}{A sufficient finite screening budget}{theorem.5}{PhysicaA_Supplement.pdf}}
\newlabel{S-eq:screen_power}{{S7}{4}{A sufficient finite screening budget}{equation.2.7}{PhysicaA_Supplement.pdf}}
\newlabel{S-prop:screen_drift}{{S6}{5}{Positive-drift screening cost}{theorem.6}{PhysicaA_Supplement.pdf}}
\newlabel{S-eq:screen_drift_cost}{{S8}{5}{Positive-drift screening cost}{equation.2.8}{PhysicaA_Supplement.pdf}}
\newlabel{S-prop:screen_break_even}{{S7}{5}{A sufficient copy-cost break-even condition}{theorem.7}{PhysicaA_Supplement.pdf}}
\newlabel{S-prop:screen_lower_bound}{{S8}{5}{An endpoint-only zero-width lower bound}{theorem.8}{PhysicaA_Supplement.pdf}}
\newlabel{S-eq:screen_lower_bound}{{S9}{5}{An endpoint-only zero-width lower bound}{equation.2.9}{PhysicaA_Supplement.pdf}}
\newlabel{S-eq:signed_screen_lower_bound}{{S10}{6}{An endpoint-only zero-width lower bound}{equation.2.10}{PhysicaA_Supplement.pdf}}
\newlabel{S-cor:implicit_region}{{S9}{6}{Implicit action regions and adaptive proposals}{theorem.9}{PhysicaA_Supplement.pdf}}
\newlabel{S-prop:race}{{S10}{7}{Stagewise approximate greedy selection}{theorem.10}{PhysicaA_Supplement.pdf}}
\newlabel{S-sec:witness_bound}{{S2.5}{7}{A fixed-witness bound for final-action certification}{subsection.2.5}{PhysicaA_Supplement.pdf}}
\newlabel{S-thm:witness}{{S11}{7}{Fixed-witness certification}{theorem.11}{PhysicaA_Supplement.pdf}}
\newlabel{S-sec:conditional_routing}{{S2.6}{8}{Independent-pilot routing and conditional risk reservation}{subsection.2.6}{PhysicaA_Supplement.pdf}}
\newlabel{S-thm:conditional_routing}{{S13}{8}{Conditional routing validity}{theorem.13}{PhysicaA_Supplement.pdf}}
\newlabel{S-sec:scoreguarantee}{{S3.1}{9}{A coefficient condition for the frozen score}{subsection.3.1}{PhysicaA_Supplement.pdf}}
\newlabel{S-eq:quadratic}{{S12}{9}{A coefficient condition for the frozen score}{equation.3.12}{PhysicaA_Supplement.pdf}}
\newlabel{S-eq:coefficients}{{S13}{9}{A coefficient condition for the frozen score}{equation.3.13}{PhysicaA_Supplement.pdf}}
\newlabel{S-thm:proxy}{{S14}{9}{Linear-exposure dominance}{theorem.14}{PhysicaA_Supplement.pdf}}
\newlabel{S-eq:dominance}{{S14}{9}{Linear-exposure dominance}{equation.3.14}{PhysicaA_Supplement.pdf}}
\newlabel{S-eq:kappabound}{{S15}{9}{Linear-exposure dominance}{equation.3.15}{PhysicaA_Supplement.pdf}}
\newlabel{S-cor:proxygreedy}{{S15}{9}{Exact greedy for the score}{theorem.15}{PhysicaA_Supplement.pdf}}
\newlabel{S-eq:proxyapprox}{{S16}{9}{Exact greedy for the score}{equation.3.16}{PhysicaA_Supplement.pdf}}
\newlabel{S-prop:localization}{{S16}{10}{Localization of positive score complementarity}{theorem.16}{PhysicaA_Supplement.pdf}}
\newlabel{S-sec:certificate}{{S3.2}{10}{Independent paired certification}{subsection.3.2}{PhysicaA_Supplement.pdf}}
\newlabel{S-eq:paired}{{S17}{10}{Independent paired certification}{equation.3.17}{PhysicaA_Supplement.pdf}}
\newlabel{S-thm:certificate}{{S17}{10}{One-query switching certificate}{theorem.17}{PhysicaA_Supplement.pdf}}
\newlabel{S-eq:eb}{{S18}{10}{One-query switching certificate}{equation.3.18}{PhysicaA_Supplement.pdf}}
\newlabel{S-eq:false}{{S19}{10}{One-query switching certificate}{equation.3.19}{PhysicaA_Supplement.pdf}}
\newlabel{S-eq:risk}{{S20}{10}{One-query switching certificate}{equation.3.20}{PhysicaA_Supplement.pdf}}
\newlabel{S-cor:portfolio}{{S18}{11}{Portfolio regret and time-uniform validity}{theorem.18}{PhysicaA_Supplement.pdf}}
\newlabel{S-eq:spending}{{S21}{11}{Portfolio regret and time-uniform validity}{equation.3.21}{PhysicaA_Supplement.pdf}}
\newlabel{S-tab:cohorts}{{S1}{11}{Evaluation cohorts, independently sampled units, and charged budgets. Final evaluation is separate from online selection}{table.1}{PhysicaA_Supplement.pdf}}
\newlabel{S-sec:conditional_details}{{S6.1}{12}{Calibration and independent confirmation}{subsection.6.1}{PhysicaA_Supplement.pdf}}
\newlabel{S-tab:robustness}{{S2}{13}{Post-hoc descriptive FEX robustness on the 60 confirmation graphs, averaging only the five eligible cells. No new inferential interval is attached}{table.2}{PhysicaA_Supplement.pdf}}
\newlabel{S-tab:grid}{{S3}{13}{Complete confirmation grid. Eligibility is fixed by independent calibration. Gain is NEX minus FEX; intervals adjust the 24 cellwise comparisons}{table.3}{PhysicaA_Supplement.pdf}}
\newlabel{S-sec:adaptive_study}{{S7}{14}{Independent-pilot routing and larger portfolios}{section.7}{PhysicaA_Supplement.pdf}}
\newlabel{S-tab:adaptive_cells}{{S4}{14}{Independent-pilot study with the frozen strong plain comparator. Every context is retained and the same exact extinction presolve is available to both methods. The common cap is 32,768 marginal copies}{table.4}{PhysicaA_Supplement.pdf}}
\newlabel{S-tab:adaptive_families}{{S5}{14}{All family-level plain/routed-screen wall ratios, reported descriptively without family-specific superiority tests}{table.5}{PhysicaA_Supplement.pdf}}
\newlabel{S-fig:adaptive_comparison}{{S1}{15}{Independent-pilot routing. (a) Mean charged copies and certified-query counts, including the common exact presolve. (b) The pooled planned plain/routed-screen geometric wall-time ratio with its conservative 97.5\% interval, and family-level descriptive points without intervals. Values above one favor routed screening. All pilot costs are included}{figure.1}{PhysicaA_Supplement.pdf}}
\newlabel{S-fig:pool_sensitivity}{{S2}{17}{Candidate-domain sensitivity, separated from computational scaling. Points show NEX-minus-FEX loss for all six size/pool settings, with the original eight-endpoint-family adjusted $99.375\%$ intervals. Distinct pool definitions are categorical and are not joined by lines. All-node pools contain $n$ candidates. Only the $n=300$, all-node interval excludes zero}{figure.2}{PhysicaA_Supplement.pdf}}
\newlabel{S-fig:complete_search_mechanisms}{{S5}{22}{Prespecified diagnostic workloads. (a) Native time on identical recorded query portfolios. (b) All candidate-count and overlap workloads; the two portfolio constructions also differ in their constituent sets, so these are not isolated causal effects of overlap}{figure.5}{PhysicaA_Supplement.pdf}}
\newlabel{S-fig:portfolio}{{S6}{23}{Portfolio evaluation at unchanged graphical joint law. (a) Graph-level geometric mean CPU-time ratios with three-comparison adjusted 98.333\% bootstrap intervals; ratios above one favor local adaptive batching. (b) Every family/size/portfolio replay-to-local ratio, shown descriptively without per-cell significance claims. Incremental RTC is a distinct coupling; its runtime comparison is at equal marginal-continuation budget}{figure.6}{PhysicaA_Supplement.pdf}}
\newlabel{S-sec:envelope_experiment}{{S18.2}{23}{Order-informed portfolio selection}{subsection.18.2}{PhysicaA_Supplement.pdf}}
\newlabel{S-fig:batch_scaling}{{S7}{24}{Descriptive portfolio scaling on fresh larger graphs. Lines show median CPU seconds per marginal continuation over eight graph instances per size; shading denotes their interquartile range, not a confidence interval. All methods evaluate the same sets and two tapes per workload. These randomized augmented states test computation and are not naturally generated histories}{figure.7}{PhysicaA_Supplement.pdf}}
\newlabel{S-tab:envelope_selection}{{S14}{25}{Fresh finite-portfolio selection: certification, endpoint-only returns, total charged online copies and wall time. Means are over all 96 contexts, including unresolved cases}{table.14}{PhysicaA_Supplement.pdf}}
\newlabel{S-fig:envelope_selection}{{S8}{25}{Order-informed selection on the fresh locked cohort. (a) Comparator/envelope-first wall-time ratios; pooled planned contrasts use adjusted 97.5\% intervals, and family intervals are descriptive 95\% intervals. Values above one favor envelope-first selection. (b) Every context's total charged online copies; auxiliary endpoints count in full. Costs include unresolved runs; neither planned comparison meets the matched-certified-quality gate. These are fixed-portfolio tolerance certificates, not guarantees of globally optimal intervention}{figure.8}{PhysicaA_Supplement.pdf}}
\newlabel{S-sec:factorial}{{S18.3}{25}{Separating endpoint information from coupling implementation}{subsection.18.3}{PhysicaA_Supplement.pdf}}
\newlabel{S-tab:factorial_cells}{{S15}{26}{Independent engine/screen study. Means include unresolved queries; endpoint-only certificates do not simulate the individual candidates}{table.15}{PhysicaA_Supplement.pdf}}
\newlabel{S-tab:factorial_cost}{{S16}{26}{All four planned graph-level geometric capped-time contrasts on the new cohort. Ratios above one favor the denominator; intervals are adjusted 98.75\% intervals. Graph signs refer to mean log ratios. Certification gates must also pass for an equal-quality efficiency claim}{table.16}{PhysicaA_Supplement.pdf}}
\newlabel{S-fig:factorial_selection}{{S9}{26}{Separating screening from the simulation engine. (a) All 96 query costs per cell; horizontal bars are medians, red crosses denote unresolved outputs, and labels give certification counts. (b) All four planned geometric wall-time ratios and adjusted 98.75\% intervals. Costs include unresolved outputs and cannot alone establish equal-quality acceleration}{figure.9}{PhysicaA_Supplement.pdf}}
\newlabel{S-sec:witness_study}{{S18.4}{27}{Independent evaluation of the fixed-witness refinement}{subsection.18.4}{PhysicaA_Supplement.pdf}}
\newlabel{S-tab:witness_cells}{{S17}{27}{Fixed-witness refinement on the independently planned third cohort. Means include every query; the common cap is 32,768 marginal copies}{table.17}{PhysicaA_Supplement.pdf}}
\newlabel{S-fig:witness_comparison}{{S10}{28}{Fixed-witness refinement. (a) Mean charged copies and certified-query counts. (b) Both planned geometric wall-time ratios with adjusted 97.5\% intervals. Values above one favor the denominator. Certification counts and costs are reported separately; the common cap exceeds the earlier studies' cap}{figure.10}{PhysicaA_Supplement.pdf}}
\newlabel{S-fig:complete_search}{{S11}{30}{Complete adaptive search. (a) The primary replay/local ratio with its $95\%$ interval and two secondary ratios with adjusted $97.5\%$ intervals. The open RTC marker denotes equal budget, rather than matched selection quality. (b) All graph-level replay/local ratios; blue and coral indicate sizes $100$ and $300$, and horizontal bars show descriptive group geometric means}{figure.11}{PhysicaA_Supplement.pdf}}
\newlabel{S-fig:optimized_replay}{{S12}{30}{Persistent-topology replay confirmation. (a) Both planned whole-search ratios with Bonferroni-adjusted $97.5\%$ intervals. (b) All 32 independent graph-level persistent-replay/local ratios, grouped by family and node count. The dashed reference denotes equal time; plotted group means are descriptive}{figure.12}{PhysicaA_Supplement.pdf}}
\newlabel{S-sec:design_details}{{S24.1}{30}{Complete history and parameter specification}{subsection.24.1}{PhysicaA_Supplement.pdf}}
\newlabel{S-sec:search_details}{{S24.2}{31}{Frozen proposal score and search specification}{subsection.24.2}{PhysicaA_Supplement.pdf}}
\newlabel{S-eq:score}{{S22}{31}{Frozen proposal score and search specification}{equation.24.22}{PhysicaA_Supplement.pdf}}
\newlabel{S-sec:portfolio_details}{{S24.3}{32}{Order-informed certificates and their role}{subsection.24.3}{PhysicaA_Supplement.pdf}}
\newlabel{S-eq:portfolio_details}{{S23}{32}{Order-informed certificates and their role}{equation.24.23}{PhysicaA_Supplement.pdf}}
\newlabel{S-sec:variance_details}{{S24.4}{32}{Fixed-contrast RTC comparison}{subsection.24.4}{PhysicaA_Supplement.pdf}}
\newlabel{S-sec:scaling_details}{{S24.5}{32}{Full scaling and pool sensitivity}{subsection.24.5}{PhysicaA_Supplement.pdf}}
\newlabel{S-fig:scalability}{{S13}{33}{Computational scaling. (a) Per-trajectory wall-time medians and interquartile ranges over sixteen measured graph instances per size, for graphical coupling and incremental RTC. Shading is an interquartile range, not a confidence interval. (b) Graphical wall time versus the measured proposal-work proxy, showing every scaling observation; the dashed descriptive log--log fit has slope $1.003$. The figure characterizes measured computational work rather than a uniform runtime advantage. Candidate-pool search effects are reported separately in Supplementary Fig.~\ref {fig:pool_sensitivity}}{figure.13}{PhysicaA_Supplement.pdf}}
\newlabel{S-sec:transfer_details}{{S24.6}{33}{Strong baselines and observed-group transfer}{subsection.24.6}{PhysicaA_Supplement.pdf}}
\newlabel{S-sec:certification_details}{{S24.7}{33}{Optional certification and charged cost}{subsection.24.7}{PhysicaA_Supplement.pdf}}
\newlabel{S-sec:quality_estimates}{{S24.8}{33}{Additional fixed-quality loss estimates}{subsection.24.8}{PhysicaA_Supplement.pdf}}
\newlabel{S-sec:memory_response_details}{{S25}{33}{Controlled-initial-memory response experiment}{section.25}{PhysicaA_Supplement.pdf}}
\newlabel{S-fig:original_memory_contexts}{{S14}{35}{Separate contexts underlying main Fig.~\ref {fig:memory_response}. Each row is one fixed graph and each column a face-strength setting. Teal and red denote initial $-T_0$ and $+T_0$; solid and dashed lines denote no control and four targets. Curves average 128 tapes within each context. The vertical line marks withdrawal. Node trajectories for the two memory initializations coincide at $s=0$}{figure.14}{PhysicaA_Supplement.pdf}}
\newlabel{S-sec:response_extension}{{S26}{35}{Memory-mechanism and comparison-bound diagnostics}{section.26}{PhysicaA_Supplement.pdf}}
\newlabel{S-tab:memory_modes}{{S20}{36}{Mean occupation loss across the eight fixed contexts. Dynamic-setting labels give initial memory and $\lambda $. All five controls are included}{table.20}{PhysicaA_Supplement.pdf}}
\newlabel{S-fig:bound_diagnostics}{{S15}{37}{All 72 empirical variance-to-bound ratios for nested and nonnested comparisons. Colors and symbols identify families. The horizontal dotted line is equality with the plug-in bound. Ratios describe the empirical distribution of tapes; they do not certify population tightness}{figure.15}{PhysicaA_Supplement.pdf}}
\newlabel{S-fig:memory_contexts}{{S16}{37}{Initial-memory loss gaps for every graph. Red: no control; teal: four targets. Solid and dashed lines distinguish the two graphs in each family. Every estimated graph-specific gap decreases over the three tested rates; no continuous-parameter monotonicity is inferred}{figure.16}{PhysicaA_Supplement.pdf}}
\newlabel{S-sec:threshold_sensitivity}{{S27}{37}{Post-hoc sensitivity to the selection-quality threshold}{section.27}{PhysicaA_Supplement.pdf}}
\newlabel{S-tab:threshold_sensitivity}{{S21}{38}{Post-hoc stable-crossing budgets in physical trajectories. Pooled means use 32 graphs; family means use eight. Thresholds and family summaries here carry no new inferential claims}{table.21}{PhysicaA_Supplement.pdf}}
\newlabel{S-sec:runtime_implementation}{{S28.1}{38}{Runtime accounting and persistent replay}{subsection.28.1}{PhysicaA_Supplement.pdf}}
\newlabel{S-sec:search_implementation}{{S28.2}{39}{Candidate screening and search}{subsection.28.2}{PhysicaA_Supplement.pdf}}
\newlabel{S-sec:natural_memory_details}{{S29}{39}{Natural-history memory and feedback study}{section.29}{PhysicaA_Supplement.pdf}}
\newlabel{S-tab:natural_memory_all}{{S22}{39}{All natural-history loss contrasts. Positive values mean the first rule has greater occupation loss. F: dynamic natural minus frozen; Z: dynamic natural minus zero; R: dynamic natural minus reset dynamic. Control 0 is absent and 1 uses four targets. Signs count graph-level means}{table.22}{PhysicaA_Supplement.pdf}}
\newlabel{S-sec:full_computational_application}{{S30}{40}{Simulation implementation and search applications}{section.30}{PhysicaA_Supplement.pdf}}
\newlabel{S-prop:work}{{S19}{40}{Work bound and exact extinction stopping}{theorem.19}{PhysicaA_Supplement.pdf}}
\newlabel{S-eq:work}{{S24}{40}{Work bound and exact extinction stopping}{equation.30.24}{PhysicaA_Supplement.pdf}}
\newlabel{S-alg:simulation}{{S2}{41}{Simulation implementation and search applications}{algorithm.S2}{PhysicaA_Supplement.pdf}}
\newlabel{S-prop:unbiased}{{S20}{41}{Fixed-query unbiasedness}{theorem.20}{PhysicaA_Supplement.pdf}}
\newlabel{S-alg:batch}{{S3}{42}{Batch evaluation of candidate actions}{algorithm.S3}{PhysicaA_Supplement.pdf}}
\newlabel{S-sec:search}{{S30.2}{42}{Search application}{subsection.30.2}{PhysicaA_Supplement.pdf}}
\newlabel{S-eq:portfolio_main}{{S25}{42}{Order-informed certification}{equation.30.25}{PhysicaA_Supplement.pdf}}
\newlabel{S-alg:search}{{S4}{43}{Anchor, score, and exchange neighborhoods}{algorithm.S4}{PhysicaA_Supplement.pdf}}
\newlabel{S-fig:cost}{{S17}{43}{Cost and selection quality for four equal-size nonnested actions. Panels (a) and (b) compare physical trajectory budgets and wall time on a shared vertical scale. Insets enlarge the original high-budget Graphical/RTC and thinning/antithetic comparisons. Markers alternate across observed budgets to reduce occlusion; all mean values and interval endpoints are reported. The reference is chosen on independent sample A and assessed on sample B. Trajectory counts charge both antithetic partners. The fixed threshold is excess loss $0.02$; curve error bars are descriptive 95\% intervals, while cost ratios use the adjusted intervals in Table~\ref {tab:cost}. Wall time covers online selection after candidates are fixed. The zero-simulation R-degree anchor is reported numerically below the panels and lies above their common vertical range}{figure.17}{PhysicaA_Supplement.pdf}}
\newlabel{S-tab:cost}{{S23}{44}{Cost at the fixed quality target: comparator cost divided by graphical cost. All intervals are Bonferroni-adjusted 99.375\% intervals for the eight-endpoint family}{table.23}{PhysicaA_Supplement.pdf}}
\newlabel{S-fig:jcs_runtime}{{S18}{45}{Complete-search computation and graph-level consistency. (a) The 64-graph study reports its primary replay/local ratio with a $95\%$ interval and secondary ratios with adjusted $97.5\%$ intervals. The open RTC marker denotes equal marginal budget and a different coupling, without matched-quality evidence. (b) The 32-graph confirmation uses persistent topology and one full-prefix call per action and stage; both intervals are adjusted $97.5\%$. (c,d) Every graph-level primary ratio is shown, distinguished by network size; short vertical bars are descriptive family geometric means. Values above one favor local batching. Primary ratios correspond to geometric-ratio time savings of 75.7\% and 73.6\%. Each cohort uses its own multiple-comparison adjustment}{figure.18}{PhysicaA_Supplement.pdf}}
\newlabel{S-fig:confirmation}{{S19}{46}{Where the conditional search gain occurs. All methods use $1,024$ online continuations per policy, with independent final evaluation. (a,b) The complete eight-cell confirmation grid on 60 graphs: entries are baseline loss minus FEX loss, averaged over histories and then within each family or across all graphs. Both panels use the same zero-centered linear color scale. Filled markers identify the five cells selected by independent calibration; open markers identify the other cells. Here $s$ is the face-strength scale and $u=3$ throughout. Family cells and pooled cell means are descriptive; the separately adjusted cellwise intervals remain in Supplementary Table~\ref {tab:grid}. (c) Aggregation over the five calibrated cells. Aggregate primary FEX--NEX and secondary FEX--greedy intervals are adjusted $97.5\%$ intervals for the two planned comparisons; family intervals are descriptive $95\%$ intervals. Histories and cells stay grouped within each graph. Positive differences favor FEX}{figure.19}{PhysicaA_Supplement.pdf}}
\newlabel{S-fig:gridheatmap}{{S19}{46}{Where the conditional search gain occurs. All methods use $1,024$ online continuations per policy, with independent final evaluation. (a,b) The complete eight-cell confirmation grid on 60 graphs: entries are baseline loss minus FEX loss, averaged over histories and then within each family or across all graphs. Both panels use the same zero-centered linear color scale. Filled markers identify the five cells selected by independent calibration; open markers identify the other cells. Here $s$ is the face-strength scale and $u=3$ throughout. Family cells and pooled cell means are descriptive; the separately adjusted cellwise intervals remain in Supplementary Table~\ref {tab:grid}. (c) Aggregation over the five calibrated cells. Aggregate primary FEX--NEX and secondary FEX--greedy intervals are adjusted $97.5\%$ intervals for the two planned comparisons; family intervals are descriptive $95\%$ intervals. Histories and cells stay grouped within each graph. Positive differences favor FEX}{figure.19}{PhysicaA_Supplement.pdf}}
\newlabel{S-tab:jcs_cohorts}{{S24}{47}{Principal evaluation cohorts. A continuation is one physical trajectory. Inference groups all histories and settings of each graph. Supplementary Table~\ref {tab:cohorts} lists all additional cohorts and budgets}{table.24}{PhysicaA_Supplement.pdf}}
\newlabel{S-sec:earlier_memory_cohorts}{{S31}{47}{Earlier controlled-memory cohorts and full design}{section.31}{PhysicaA_Supplement.pdf}}
\newlabel{S-sec:verification}{{S31.2}{48}{Exactness and mathematical verification}{subsection.31.2}{PhysicaA_Supplement.pdf}}
\newlabel{S-sec:memory_response}{{S31.3}{48}{Controlled initial memory and spreading response}{subsection.31.3}{PhysicaA_Supplement.pdf}}
\newlabel{S-sec:timescales}{{S31.4}{48}{Static bias, dynamic memory and time scales}{subsection.31.4}{PhysicaA_Supplement.pdf}}
\newlabel{S-fig:memory_response}{{S20}{49}{Controlled-initial-memory response (descriptive). Columns show $s=0$ and $s=1$; rows show mean rumor occupation and mean face memory. Curves average $128$ trajectories on each of four fixed graph/initial-state contexts. Color denotes initial face memory; line style denotes intervention. Shading shows conditional Monte Carlo standard errors, not uncertainty over a graph population. The vertical line marks control withdrawal at $D=2$. The two memory-initialization node curves coincide when $s=0$. All conditions are shown}{figure.20}{PhysicaA_Supplement.pdf}}
\newlabel{S-sec:natural_memory}{{S31.5}{49}{Naturally generated memories and feedback}{subsection.31.5}{PhysicaA_Supplement.pdf}}
\newlabel{S-fig:timescales}{{S21}{50}{Memory mechanisms at fixed higher-order transmission (descriptive; eight graphs). Red denotes no control and teal denotes four targets. (a) Loss difference between initial $+T_0$ and $-T_0$ under dynamic memory; dotted lines show frozen-memory differences. Error bars are one conditional Monte Carlo standard error for the equally weighted graph average; they quantify tape noise with these graphs fixed, not uncertainty across a population of graphs. (b) Mean loss from initial $-T_0$ (solid circles) and $+T_0$ (dashed squares); dotted horizontal lines show memory fixed at zero. Both signs use identical initial node states. Frozen memory means $\lambda =0$; the other tick labels give the relaxation rate. Lines connect tested settings, with no interpolation-based inference}{figure.21}{PhysicaA_Supplement.pdf}}
\newlabel{S-fig:natural_memory}{{S22}{51}{Naturally generated memory (32 graphs, 64 histories; descriptive). (a) Post-control occupation under four memory rules, with four targets. Frozen and zero controls are identical across the continuation-rate grid. (b) Spatial mean-memory contrast between dynamic natural and dynamic reset branches after subtracting direct exponential decay of the initial memory; the vertical line marks control withdrawal. Full pointwise uncertainty and both control conditions are reported in the supplement and accompanying CSV files}{figure.22}{PhysicaA_Supplement.pdf}}
\newlabel{S-sec:physics_new}{{S32}{51}{Withdrawal-time experiments and observed contact data}{section.32}{PhysicaA_Supplement.pdf}}
\newlabel{S-tab:withdrawal_all}{{S25}{52}{Complete withdrawal-time results. $G$ is total loss reduction, $G_A$ its memory component, and $\eta _A$ the memory fraction. The interval applies to $\eta _A$ in percent}{table.25}{PhysicaA_Supplement.pdf}}
\newlabel{S-sec:triangle_new}{{S33}{53}{Single-triangle retention and rebound}{section.33}{PhysicaA_Supplement.pdf}}
\newlabel{S-tab:triangle_all}{{S26}{53}{All triangle cells. $\mathcal R$ is mean sampled pathwise rebound; $\mathcal R_{\rm mean}$ is the positive overshoot of the mean curve; memory fractions are in percent}{table.26}{PhysicaA_Supplement.pdf}}
\newlabel{S-fig:triangle_ensemble}{{S23}{54}{Overshoot of the mean occupation curve on the same complete triangle grid. Its scale is much smaller than mean within-path rebound}{figure.23}{PhysicaA_Supplement.pdf}}
\newlabel{S-sec:hyperedge4}{{S34}{54}{Weighted four-uniform factors}{section.34}{PhysicaA_Supplement.pdf}}
\newlabel{S-fig:hyperedge4}{{S24}{55}{Ordered occupation response on weighted four-uniform factors under two bounded antitone drives. The vertical line marks withdrawal}{figure.24}{PhysicaA_Supplement.pdf}}
\newlabel{S-sec:precision_new}{{S35}{55}{Threshold-free precision and cost}{section.35}{PhysicaA_Supplement.pdf}}
